\chapter{Decadimenti radioattivi}
\section{Reazioni nucleari}
Nei processi di decadimento radioattivo i nuclei decadono spontaneamente emettendo particelle e/o radiazioni elettromagnetiche, trasformandosi in nuclei di altri elementi. Alternativamente possono essere indotti a decadere tramite l'interazione con particelle esterne, ad esempio bombardando il nucleo con neutroni o protoni. Queste sono le \textbf{reazioni nucleari}:
\begin{equation}
    a + X \to Y + b
\end{equation}
Dove $a$ e $b$ sono particelle (ad esempio neutroni, protoni, particelle $\alpha$, ecc.) e $X$ e $Y$ sono i nuclei coinvolti nella reazione. Durante una reazione nucleare l'energia di reazione è data da:
\begin{equation}
    Q = (M_X + M_a - M_Y - M_b)c^2
\end{equation}
Se $Q > 0$ la reazione è esotermica (rilascia energia) dove osserviamo un aumento dell'energia cinetica degli elementi coinvolti, se $Q < 0$ la reazione è endotermica (assorbe energia), possibile solo se l'energia cinetica di $a$ è maggiore di $\abs{Q}$.\\
Oltre ad energia e quantità di moto, durante le reazioni nucleari si conserva anche il numero totale di nucleoni.\\
Un nuclide radioattivo decade spontaneamente in liberando una o più particelle o radiazioni elettromagnetiche. Possiamo separare le particelle emesse grazie ad un campo magnetico, ad esempio:
\begin{itemize}
    \item Particelle $\alpha$: nuclei di elio composti da 2 protoni e 2 neutroni, con carica $+2e$
    \item Elettroni ad alta energia (particelle $\beta^-$) con carica $-e$
    \item Raggi $\gamma$: radiazione elettromagnetica ad alta energia, senza carica
\end{itemize}
Il decadimento radioattivo è un processo stocastico, ovvero non è possibile prevedere quando un singolo nucleo decadrà, ma è possibile prevedere il comportamento di un insieme di nuclei. La quantità di nuclei $N$ che decadono in un intervallo di tempo $\dd{t}$ è proporzionale al numero di nuclei presenti e all'intervallo di tempo:
\begin{equation}
    \dv{\mathcal{N}}{t} = -\lambda \mathcal{N} \implies \mathcal{N}(t) = \mathcal{N}_0 \efunction^{-\lambda t}
\end{equation}
Dove $\lambda$ è la costante di decadimento, caratteristica di ogni nuclide.
Possiamo anche definire l'\textbf{attività} $\mathcal{A}$ come il numero di decadimenti per unità di tempo:
\begin{equation}
    \mathcal{A}(t) = \lambda \mathcal{N}(t) = \lambda \mathcal{N}(t) = \mathcal{A}_0 \efunction^{-\lambda t}
\end{equation}
La vita media di un nuclide è data da:
\begin{equation}
    \tau = \avg{\mathcal{A}(t)} = \frac{1}{\lambda}
\end{equation}
L'unità di misura dell'attività nel SI è il becquerel (Bq), che corrisponde ad un decadimento al secondo. Un'altra unità di misura comune è la curie (Ci), definita come $1 \text{ Ci} = \qty{3.7e10}{\becquerel}$.\\
Infine definiamo il \textbf{tempo di dimezzamento} $T_{1/2}$ come l'intervallo di tempo necessario affinché la quantità di nuclei si riduca alla metà del valore iniziale:
\begin{equation}
    \mathcal{N}(t + T_{1/2}) = \dfrac{\mathcal{N}(t)}{2} \implies \dfrac{1}{2}\cancel{\mathcal{N}_0 \efunction^{-\lambda t}} = \cancel{\mathcal{N}_0 \efunction^{-\lambda t}}\efunction^{-\lambda T_{1/2}} \implies T_{1/2} = \frac{\ln 2}{\lambda}
\end{equation}

\addsvg[Decadimento radioattivo e tempo di dimezzamento]{Radioactive_decay}{0.6}

\section{Decadimento alpha}
Il decadimento $\alpha$ avviene quando un nucleo instabile emette una particella $\alpha$ (nucleo di elio) trasformandosi in un nucleo di un elemento con numero atomico diminuito di 2 e numero di massa diminuito di 4:
\begin{equation}
    \atom{X}{A}{Z} \to \atom{Y}{A-4}{Z-2} + \atom{He}{4}{2} = \atom{Y}{A-4}{Z-2} + \alpha
\end{equation}
Ad esempio il Radium-226 decade in Radon-222 emettendo una particella $\alpha$:
\begin{equation}
    \atom{Ra}{226}{88} \to \atom{Rn}{222}{86} + \alpha
\end{equation}
Se ci poniamo nel sistema di riferimento del nucleo padre, la particella $\alpha$ e il nucleo figlio si allontanano con velocità opposte per conservare la quantità di moto. Indicando con $K_\alpha$ l'energia cinetica della particella $\alpha$ e con $K_Y$ l'energia cinetica del nucleo figlio, possiamo scrivere:
\begin{equation}
    Q = K_\alpha + K_Y
\end{equation}
Tipicamente $K_Y \ll K_\alpha$ poiché la massa del nucleo figlio è molto maggiore di quella della particella $\alpha$. Inoltre $K_\alpha \approx \qty{5}{\mega\electronvolt}$. Ad esempio, nel decadimento dell'uranio-238 in torio-234, l'energia di reazione è $Q = \qty{4.27}{\mega\electronvolt}$, di cui $\qty{4.25}{\mega\electronvolt}$ sono portati dalla particella $\alpha$. La reazione è la seguente:
\begin{equation}
    \atom{U}{238}{92} \to \atom{Th}{234}{90} + \alpha
\end{equation}
Caratterizzata da un tempo di dimezzamento di $T_{1/2} = 4.47 \times 10^9 \text{ anni}$.\\
Poiché l'uranio è particolarmente pesante, possiamo distinguere una lunga catena di decadimenti $\alpha$ e $\beta$ che porta alla formazione di un nuclei più stabili, come il piombo-206:

\addfigure[Catena di decadimento dell'uranio-238]{drawio-pdf/U_238_decay.drawio.pdf}{0.7}

Per capire il funzionamento del decadimento $\alpha$, consideriamo il nucleo come un pozzo di potenziale finito, dove la particella $\alpha$ è confinata all'interno del nucleo da una barriera di potenziale, che va da circa $-\qty{40}{\mega\electronvolt}$ a $\qty{30}{\mega\electronvolt}$, La particella $\alpha$ ha un'energia cinetica di circa $\qty{5}{\mega\electronvolt}$, quindi non ha abbastanza energia per superare la barriera di potenziale e uscire dal nucleo. Tuttavia, grazie al fenomeno del tunneling quantistico, c'è una probabilità finita che la particella $\alpha$ possa attraversare la barriera di potenziale ed essere emessa dal nucleo.\\
Lo schema rappresentativo è il seguente:

\addtikz[Decadimento $\alpha$ come fenomeno di tunneling quantistico]{
    %Axis
    \draw[-stealth,thick] (-0.5,0) -- (6,0) node[right] {$r$};
    \draw[-stealth,thick] (0,-5) -- (0,4) node[above] {$U(r)$};
    %Annotations
    \draw[dashed,thick] (0,0.75) node[left]{$\qty{5}{\mega\electronvolt}$} -- (2.63,0.75) -- (2.63,0)node[below]{$R_1$};
    \draw[dashed,thick] (0,3) node[left]{$\approx \qty{30}{\mega\electronvolt}$} -- (1.5,3);
    \node[left] at (0,-4) {$\approx -\qty{40}{\mega\electronvolt}$};
    \node[below right] at (0.75,0) {$R$};
    \node[below left] at (0,0) {0};
    \node[above] at (4,3) {\centering Repulsive};
    \node[below] at (4,3) {\centering Coulomb potential};
    \draw[-stealth,thick] (4,2.25) to[out=-90,in=0] (2.5,1.5);
    \node[above] at (-1.5,-1) {\centering Attractive nuclear};
    \node[below] at (-1.5,-1) {\centering potential hole};
    \draw[-stealth,thick] (-1.5,-1.75) to[out=-90,in=180] (0.25,-2.5);
    %Potential barrier
    \draw[thick,supportDarkBlue] (0,-4) to[out=-15,in=180,looseness=0.5] (1.5,3) to[out=0,in=115,looseness=0.6] (2.25,1.5) to[out=-70,in=178,looseness=1.3] (4.5,0.1);
    \fill[black] (0.77,0.75) circle (0.05)node[above right]{A};
    \fill[black] (2.62,0.75) circle (0.05)node[above right]{B};
}{1}{tunneling_alpha_decay}

Se approssimiamo la barriera di potenziale come una barriera rettangolare di altezza $V_0$ e larghezza $L$, la probabilità di trasmissione della particella $\alpha$ attraverso la barriera è data da:
\begin{equation}
    T \approx \efunction^{-2\gamma L} \quad \text{con} \quad \gamma = \frac{\sqrt{2m(V_0 - E)}}{\hbar}
\end{equation}
Dove $m$ è la massa della particella $\alpha$, $E$ è la sua energia cinetica, e $\hbar$ è la costante di Planck ridotta. La probabilità di trasmissione dipende esponenzialmente dalla larghezza e dall'altezza della barriera di potenziale, il che spiega perché il decadimento $\alpha$ è un processo raro e perché i nuclei con barriere di potenziale più alte e larghe hanno tempi di dimezzamento più lunghi.\\
Ad esempio con una barriera di potenziale di altezza $V_0 = \qty{40}{\electronvolt}$ e larghezza $L = \qty{1}{\nano\meter}$, la probabilità di trasmissione per un elettrone $e^-$ con energia cinetica $E = \qty{30}{\electronvolt}$ è pari a:
\begin{equation}
    \begin{split}
        & 2\gamma L = 2\dfrac{\sqrt{2 \cdot \qty{9.11e-31}{\kilo\gram} \cdot (\qty{40}{\electronvolt} - \qty{30}{\electronvolt}) \cdot \qty{1.6e-19}{\joule\per\electronvolt}}}{\qty{1.05e-34}{\joule\second}} \cdot \qty{1e-9}{\meter} \approx 32.5 \\[8pt]
        & T \approx \efunction^{-32.5} \approx 8.5 \times 10^{-15} \rightarrow \text{Probabilità} \approx 10^{-12} \%
    \end{split}
\end{equation}
Se diminuiamo la larghezza della barriera a $L = \qty{0.1}{\nano\meter}$, la probabilità di trasmissione diventa:
\begin{equation}
    T \approx \efunction^{-3.25} \approx 3.9 \times 10^{-2} \rightarrow \text{Probabilità} \approx 4 \%
\end{equation}
Il modello a barriera rettangolare ci serve per ottenere la legge di Geiger-Nuttall, che mette in relazione il tempo di dimezzamento $T_{1/2}$ di un nuclide con l'energia cinetica $K_\alpha$ della particella $\alpha$ emessa. Inizialmente questa legge fu ottenuta empiricamente, ma può essere giustificata tramite il modello di tunneling quantistico. Una particella $\alpha$ che attraversa la barriera di potenziale Coulombiano dal punto $A$ al punto $B$ può essere descritta come la somma di passaggi per barriere di potenziale rettangolari di larghezza infinitesima $\Delta r_i$. La probabilità di trasmissione attraverso la barriera complessiva è data dal prodotto delle probabilità di trasmissione attraverso ogni barriera infinitesimale:
\begin{equation}
    T = \prod_i \efunction^{-2\gamma_i \Delta r_i} = \efunction^{-2\sum_i \gamma_i \Delta r_i} \xrightarrow[\Delta r_i \to 0]{} \efunction^{-2\int_A^B \gamma(r) \dd{r}}
\end{equation}
Possiamo definire il \textbf{fattore di Gamow} $G$ come:
\begin{equation}
    G = \int_R^{R_c} \gamma(r) \dd{r} = \frac{1}{\hbar} \int_R^{R_c} \sqrt{2M_{\alpha}(U_c(r) - K_\alpha)} \dd{r}
\end{equation}
Dove $R$ è il raggio del nucleo, $R_c$ è il punto in cui la barriera di potenziale Coulombiano incontra l'energia cinetica della particella $\alpha$, $M_\alpha$ è la massa della particella $\alpha$, e $U_c(r)$ è il potenziale Coulombiano. Poiché:
\begin{equation}
    \dfrac{U_c(r)}{K_\alpha} = \dfrac{R_c}{r}
\end{equation}
Riscriviamo il fattore di Gamow come:
\begin{equation}
    G = \frac{\sqrt{2M_\alpha K_\alpha}}{\hbar} \int_R^{R_c} \sqrt{\dfrac{R_c}{r} - 1} \dd{r}
\end{equation}
Effettuando il cambio di variabile $x = r/R_c$, otteniamo:
\begin{equation}
    G = \frac{R_c}{\hbar}\sqrt{2M_\alpha K_\alpha}  \int_{R/R_c}^{1} \sqrt{\dfrac{1}{x} - 1} \dd{x}
\end{equation}
Inoltre $R \ll R_c$ quindi possiamo approssimare l'integrale come:
\begin{equation}
    \int_{0}^{1} \sqrt{\dfrac{1}{x} - 1} \dd{x} = \frac{\pi}{2}
\end{equation}
Quindi il fattore di Gamow diventa:
\begin{equation}
    G \approx \frac{\pi R_c}{2\hbar} \sqrt{2M_\alpha K_\alpha}
\end{equation}
Ricordando che:
\begin{equation}
    R_c = \dfrac{1}{4\pi \varepsilon_0} \dfrac{2(Z-2) e^2}{K_\alpha}, \qquad \alpha_c = \dfrac{e^2}{4\pi \varepsilon_0 \hbar c} \approx \dfrac{1}{137}, \qquad K_\alpha = \dfrac{1}{2}M_\alpha v^2
\end{equation}
Troviamo che:
\begin{equation}
    G \approx \dfrac{\pi (Z-2) e^2}{4\pi \varepsilon_0 \hbar} \sqrt{\dfrac{2M_\alpha}{K_\alpha}} = \dfrac{(Z-2) \alpha_c \pi}{v/c}
\end{equation}
Infine, troviamo la legge di Geiger-Nuttall:
\begin{equation}
    T = \efunction^{-2G} \propto \lambda \propto \dfrac{1}{T_{1/2}} \implies \boxed{\log_{10} T_{1/2} \propto \dfrac{Z}{\sqrt{K_\alpha}}}
\end{equation}
Un esempio applicativi della legge di Geiger-Nuttall e della teoria del tunneling quantistico nel decadimento $\alpha$ sono i generatori \textbf{RTG} (Radioisotope Thermoelectric Generators), utilizzati nelle sonde spaziali per generare energia elettrica. Questi motori sfruttano il calore generato dal decadimento radioattivo di isotopi come il plutonio-238, che decade emettendo particelle $\alpha$. Il calore prodotto viene convertito in energia elettrica tramite termocoppie, fornendo una fonte di energia affidabile e duratura per le missioni spaziali a lungo termine. Lo schema di un RTG è il seguente:

\addfigure[Schema di un generatore RTG. Fonte: \href{https://en.wikipedia.org/wiki/Radioisotope_thermoelectric_generator}{Wikipedia}]{png/Cutdrawing_of_an_GPHS-RTG}{0.7}

Esistono diversi metodi di produzione del plutonio-238, riportiamo due reazioni nucleari comuni:

\addfigure[Produzione del plutonio-238 tramite reazioni nucleari]{drawio-pdf/Polonium_production.drawio.pdf}{0.8}

\section{Decadimento beta}
I decadimenti $\beta$ avvengono quando un nucleo instabile emette una particella $\beta$ (elettrone o positrone) trasformandosi in un nucleo di un elemento con numero atomico aumentato o diminuito di 1, mentre il numero di massa rimane invariato. Esempi di decadimenti $\beta$ sono:
\begin{equation}
    \begin{split}
        &\atom{P}{32}{} \to \atom{S}{32}{} + e^- + \anti{\nu}_{e} \qquad (T_{1/2} = 14.3\,\text{days}) \\[8pt]
        &\atom{Cu}{64}{} \to \atom{Ni}{64}{} + e^+ + {\nu}_{e} \qquad (T_{1/2} = 12.7\,\text{hrs})
    \end{split}
\end{equation}
Le due reazioni fondamentali che danno luogo ai decadimenti elencati sono:
\begin{equation}
    \begin{split}
        & n \to p + e^- + \anti{\nu}_{e} \quad \text{(decadimento $\beta^-$)} \\[8pt]
        & p \to n + e^+ + {\nu}_{e} \quad \text{(decadimento $\beta^+$)}
    \end{split}
\end{equation}
La caratteristica fondamentale dei decadimenti $\beta$ è che l'energia cinetica delle particelle emesse non è discreta, ma segue uno spettro continuo. Questo fenomeno fu inizialmente inspiegabile, poiché si pensava che l'energia cinetica dell'elettrone dovesse essere fissa in base alla differenza di massa tra il nucleo padre e il nucleo figlio, ovvero con la stessa reazione nucleare ottengo la stessa energia per la particella $\beta$. Tuttavia, l'ipotesi dell'esistenza del neutrino da parte di Wolfgang Pauli nel 1930 risolse questo enigma. Il neutrino è una particella neutra e quasi priva di massa che viene emessa insieme all'elettrone durante il decadimento $\beta$. La presenza del neutrino consente di conservare l'energia e la quantità di moto nel processo di decadimento, spiegando lo spettro continuo dell'energia cinetica dell'elettrone.\\
La scoperta sperimentale del neutrino avvenne nel 1956 grazie agli esperimenti di Clyde Cowan e Frederick Reines, che osservarono le interazioni dei neutrini prodotti da un reattore nucleare.\\
Per studiare i neutrini possiamo osservare i prodotti delle reazioni all'interno di reattori a fissione:
\begin{equation}
    n \to p + e^- + \anti{\nu}_{e}
\end{equation}
Oppure realizzare rilevatori, costituiti da una vasca di 200 litri di acqua mescolata a cloruro di sodio e cadmio, circondata da rivelatori di scintillazione. I neutrini interagiscono con i protoni nell'acqua, producendo positroni che generano segnali nei rivelatori:
\begin{equation}
    \anti{\nu}_{e} + p \to n + e^+
\end{equation}
Ricordando che la produzione di positroni è seguita dall'annichilazione con elettroni, producendo raggi gamma:
\begin{equation}
    e^+ + e^- \to 2\gamma_{e}
\end{equation}
E sapendo che il neutrone prodotto viene catturato dal cadmio, emettendo un altro segnale gamma:
\begin{equation}
    n + \atom{Cd}{113}{} \to \atom{Cd}{114}{} + \gamma_n
\end{equation}
Possiamo rilevare i neutrini osservando le coincidenze tra i segnali prodotti dai positroni e dai neutroni. Questo metodo ha permesso di confermare l'esistenza dei neutrini e di studiarne le proprietà fondamentali. In particolare abbiamo che $\gamma_n \approx \qty{9}{\mega\electronvolt}$ e $\gamma_{e} \approx \qty{0.511}{\mega\electronvolt}$.\\
Inoltre possiamo confermare che i neutrini non sono uguali alle loro antiparticelle poichè, sperimentalmente non osserviamo mai il processo inverso:
\begin{equation}
    {\nu}_{e} + n \not\to p + e^-
\end{equation}
Questo perché i neutrini interagiscono in modo estremamente debole con la materia, rendendo molto improbabile l'osservazione di tale processo. Per interagire un $\nu_{e}$ dovrebbe attraversare uno stato di materia di lunghezza circa pari a $10\,\text{ly} \approx \qty{e13}{\kilo\meter}$ di piombo.\\
La generica reazione di decadimento $\beta$ può essere rappresentata come:
\begin{equation}
    \atom{Y}{A}{Z} \to \atom{X}{A}{Z \pm 1} + e^{\mp} + (\anti{\nu}_{e} \text{ o } {\nu}_{e})
\end{equation}
Dove il segno $+$ si riferisce al decadimento $\beta^-$ e il segno $-$ al decadimento $\beta^+$.\\
L'energia di reazione $Q$ in un decadimento $\beta$ è data da:
\begin{equation}
    Q = (M_Y - M_X - m_e)c^2
\end{equation}
Ovvero:
\begin{equation}
    \begin{split}
        Q_{\beta^-} &= \sqbr{M_{\text{nuc}}(\atom{Y}{A}{Z}) - M_{\text{nuc}}(\atom{X}{A}{Z+1}) - m_e}c^2 =\\[8pt]
        &\approx  \sqbr{M_{\text{at}}(\atom{Y}{A}{Z}) - Zm_e}c^2 - \sqbr{M_{\text{at}}(\atom{X}{A}{Z+1}) - (Z+1)m_e}c^2 - m_e c^2 =\\[8pt]
        &= \sqbr{M_{\text{at}}(\atom{Y}{A}{Z}) - M_{\text{at}}(\atom{X}{A}{Z+1})}c^2
    \end{split}
\end{equation}
\begin{equation}
    \begin{split}
        Q_{\beta^+} &= \sqbr{M_{\text{nuc}}(\atom{Y}{A}{Z}) - M_{\text{nuc}}(\atom{X}{A}{Z-1}) - m_e}c^2 =\\[8pt]
        &\approx  \sqbr{M_{\text{at}}(\atom{Y}{A}{Z}) - Zm_e}c^2 - \sqbr{M_{\text{at}}(\atom{X}{A}{Z-1}) - (Z-1)m_e}c^2 - m_e c^2 =\\[8pt]
        &= \sqbr{M_{\text{at}}(\atom{Y}{A}{Z}) - M_{\text{at}}(\atom{X}{A}{Z-1}) - 2m_e}c^2
    \end{split}
\end{equation}
Mentre nel caso della reazione inversa di cattura elettronica:
\begin{equation}
    e^- + \atom{X}{A}{Z} \to \atom{Y}{A}{Z-1} + \nu_{e}
\end{equation}
L'energia di reazione è:
\begin{equation}
    \begin{split}
        Q_{\epsilon} &= \sqbr{M_{\text{nuc}}(\atom{Y}{A}{Z}) + m_e - M_{\text{nuc}}(\atom{X}{A}{Z-1}) }c^2 =\\[8pt]
        &\approx  \sqbr{M_{\text{at}}(\atom{Y}{A}{Z}) - Zm_e}c^2 - \sqbr{M_{\text{at}}(\atom{X}{A}{Z-1}) - (Z-1)m_e}c^2 + m_e c^2 =\\[8pt]
        &= \sqbr{M_{\text{at}}(\atom{Y}{A}{Z}) - M_{\text{at}}(\atom{X}{A}{Z-1})}c^2
    \end{split}
\end{equation}
Ovvero la stessa espressione del decadimento $\beta^-$ (anche se è la reazione derivata dal decadimento $\beta^+$).\\
Un aspetto interessante dei decadimenti $\beta$ è che la parità non è conservata in questi processi. Questo fenomeno fu scoperto negli anni '50 da Chien-Shiung Wu, che osservò che gli elettroni emessi nel decadimento $\beta$ del cobalto-60 tendevano ad essere emessi in una direzione preferenziale rispetto allo spin del nucleo, violando la simmetria di parità. Questa scoperta portò alla comprensione che le interazioni deboli, responsabili dei decadimenti $\beta$, non rispettano la simmetria di parità, un risultato fondamentale nella fisica delle particelle.

\addfigure[Decadimento $\beta$ del cobalto-60 e violazione della parità]{drawio-pdf/Parity_transformations.drawio.pdf}{0.8}

In particolare osserviamo che nelle interazioni deboli l'elicità delle particelle gioca un ruolo cruciale. Gli elettroni emessi nel decadimento $\beta^-$ sono prevalentemente di elicità sinistra, mentre i positroni emessi nel decadimento $\beta^+$ sono di elicità destra.\\
Ma cosa implicava questa scoperta? Di fatto la scoperta di Wu dimostrava che le interazioni deboli hanno una preferenza verso una certa configurazione di spin e momento, violando la simmetria di parità. In particolare si osserva che gli elettroni hanno preferibilmente elicità negativa (sinistra) e i neutrini hanno preferibilmente elicità positiva (destra). Successivamente si scoprirà che i neutrini hanno sempre elicità positiva, mentre gli antineutrini hanno sempre elicità negativa.
\section{Decadimenti gamma}
I decadimenti $\gamma$ avvengono quando un nucleo eccitato rilascia energia sotto forma di radiazione elettromagnetica ad alta energia (raggi $\gamma$) per tornare al suo stato fondamentale. Questo processo non comporta un cambiamento nel numero di massa o nel numero atomico del nucleo, ma solo una transizione tra stati energetici diversi. Un esempio di decadimento $\gamma$ è:
\begin{equation}
    \atom{C}{12}{6}^* \to \atom{C}{12}{6} + \gamma
\end{equation}
Dove $\atom{C}{12}{6}^*$ rappresenta un nucleo di carbonio eccitato ad un'energia circa $\qty{4.4}{\mega\electronvolt}$ superiore a quella del suo stato fondamentale.

\addtikz[Decadimento $\gamma$ di un nucleo eccitato]{
    \draw[fermion=red] (0.5,2) -- (2.5,0)node[midway,above right] {$e^-$};
    \draw[fermion=red] (-0.5,2) -- (1.5,-1)node[midway,above right] {$e^-$};
    \node[above] at (3,0){$\atom{C}{12}{6}^*$};
    \node[above] at (3,-1){$\atom{C}{12}{6}$};
    \node[above] at (0,2){$\atom{B}{12}{5}$};
    \draw[photon=betterYellow] (2,0) -- (2,-1) node[midway,right] {$\gamma$};
    \draw[thick] (4,0) -- (5,0)
    (3.5,-1) -- (5,-1) (3.5,2) -- (5,2);
    \draw[thick,supportDarkBlue] (-1,2) -- (1,2);
    \draw[thick,supportDarkBlue] (1,0) -- (3,0);
    \draw[thick,supportDarkBlue] (1,-1) -- (3,-1);
    \draw[stealth-stealth,thick] (4.75,-1) -- (4.75,0) node[midway,right] {$\qty{4.4}{\mega\electronvolt}$};
    \draw[stealth-stealth,thick] (3.75,-1) -- (3.75,2) node[midway,right] {$\qty{13.4}{\mega\electronvolt}$};
}{1}{gamma_decay}

\section{Leptoni}
Parliamo ora dell'ultima famiglia fondamentale di particelle: i leptoni. I leptoni sono particelle elementari che non partecipano alle interazioni forti, ma sono soggetti alle interazioni elettromagnetiche, deboli e gravitazionali. I leptoni più noti sono l'elettrone ($e^-$), il muone ($\mu^-$) e il tauone ($\tau^-$), insieme ai loro rispettivi neutrini ($\nu_{e}$, $\nu_{\mu}$, $\nu_{\tau}$).\\
Ad ogni leptone carico è associato un neutrino neutro, e ogni particella ha la sua antiparticella con carica opposta. I leptoni sono fermioni, il che significa che obbediscono al principio di esclusione di Pauli e hanno spin semi-intero pari a $1/2$.\\
Il numero leptonico è una quantità conservata. Ogni leptone ha un numero leptonico di $+1$, mentre ogni antileptone ha un numero leptonico di $-1$. I neutrini e gli antineutrini hanno numeri leptonici di $+1$ e $-1$ rispettivamente. Durante le reazioni nucleari e i decadimenti, il numero totale di leptoni deve rimanere costante. Ad esempio, nel decadimento $\beta^-$:
\begin{equation}
    n \to p + e^- + \anti{\nu}_{e}
\end{equation}
Il numero leptonico totale prima del decadimento è $0$ (il neutrone non è un leptone), e dopo il decadimento è $+1$ (per l'elettrone) e $-1$ (per l'antineutrino), che si sommano a $0$, conservando il numero leptonico. Le proprietà fondamentali dei leptoni sono riassunte nella seguente tabella:
\begin{table}[H]
    \centering
    \begin{tabular}{cccccccc}
        \hline \\[-4pt]
        Simbolo & M & Q & $L_e$ & $L_\mu$ & $L_\tau$ & Tempo di vita & Decadimenti principali \\[8pt]
        \hline \\[-4pt]
        $e^-$ & $\qty{0.511}{\mega\electronvolt/c^2}$ & $-1$ & $1$ & $0$ & $0$ & Stabile & - \\[8pt]
        $\mu^-$ & $\qty{105.7}{\mega\electronvolt/c^2}$ & $-1$ & $0$ & $1$ & $0$ & $\qty{2.197e-6}{\second}$ & $e^-,\anti{\nu}_{e},\nu_{\mu}\,(100\%)$ \\[8pt]
        $\tau^-$ & $\qty{1777.0}{\mega\electronvolt/c^2}$ & $-1$ & $0$ & $0$ & $1$ & $\qty{2.906e-13}{\second}$ & $\mu^-,\anti{\nu}_{\mu},\nu_{\tau}\,(17.4\%)$ \\[8pt]
        &&&&&&& $e^-,\anti{\nu}_{e},\nu_{\tau}\,(17.8\%)$ \\[8pt]
        &&&&&&& $\nu_\tau + \text{adroni}\,(\sim 64\%)$ \\[8pt]
        $\nu_{e}$ & $< \qty{2.2}{\electronvolt/c^2}$ & $0$ & $1$ & $0$ & $0$ & Stabile & - \\[8pt]
        $\nu_{\mu}$ & $< \qty{0.19}{\mega\electronvolt/c^2}$ & $0$ & $0$ & $1$ & $0$ & Stabile & - \\[8pt]
        $\nu_{\tau}$ & $< \qty{18.2}{\mega\electronvolt/c^2}$ & $0$ & $0$ & $0$ & $1$ & Stabile & - \\[8pt]
        \hline
    \end{tabular}
    \caption{Proprietà fondamentali dei leptoni}
\end{table}
\section{Applicazioni del decadimento radioattivo}
Un'importante applicazione del decadimento $\beta^-$ è la datazione al carbonio-14, un metodo utilizzato in archeologia e geologia per determinare l'età di reperti organici. Il carbonio-14 ($\atom{C}{14}{}$) è un isotopo radioattivo del carbonio che si forma nell'atmosfera terrestre attraverso l'interazione dei raggi cosmici con l'azoto:
\begin{equation}
    \begin{split}
        &\atom{N}{14}{7} + n \to \atom{C}{14}{6} + p \\[8pt]
        &\atom{C}{14}{6} \to \atom{N}{14}{7} + e^- + \anti{\nu}_{e}
    \end{split}
\end{equation}
Una volta formato, il carbonio-14 viene assorbito dagli organismi viventi attraverso il ciclo del carbonio. La concentrazione di carbonio-14 in atmosfera rimane relativamente costante:
\begin{equation}
    \dfrac{N(\atom{C}{14}{})}{N(\atom{C}{12}{})} \approx 1.3 \times 10^{-12}
\end{equation}
Negli organismi biologici c'è un continuo riscambio di carbonio con l'ambiente, mantenendo un rapporto costante tra carbonio-14 e carbonio-12. Tuttavia, quando un organismo muore, smette di assorbire carbonio, e il carbonio-14 presente inizia a decadere con un tempo di dimezzamento di circa $T_{1/2} = 5730$ anni. Misurando la quantità di carbonio-14 rimanente in un campione organico e confrontandola con il rapporto iniziale, è possibile calcolare il tempo trascorso dalla morte dell'organismo, fornendo così una stima dell'età del reperto fino a circa 25.000 anni (per datazioni radiometriche più antiche si utilizzano altri isotopi come $\atom{K}{40}{}, \atom{U}{235}{}, \atom{Pb}{207}{}, \dots$). %da vedere slide 23 a 26
